\documentclass[aspectratio=169,10pt,spanish]{beamer}

\input{../../../_preambulo_beamer.tex}
\input{../../../_comandos_beamer.tex}

\selectlanguage{spanish}

\title{MA1001B - Modelación Estadística para la Toma de Decisiones}
\subtitle{Lección 43: Diseño experimental: factores y aleatorización}
\author{}
\institute{Tecnol\'ogico de Monterrey}
\date{}

\begin{document}

\begin{frame}[plain]
  \vspace{1.2cm}
  {\Large\bfseries MA1001B - Modelación Estadística para la Toma de Decisiones\par}
  \vspace{0.7cm}
  {\large Lección 43: Diseño experimental: factores y aleatorización\par}
  \vspace{0.7cm}
  {\color{TecRojo}\rule{\textwidth}{1.2pt}}
  \vspace{0.8cm}

  Facultad MA1001B\\[0.3cm]
  Periodo\\[0.3cm]
  Tecnol\'ogico de Monterrey
\end{frame}

\begin{frame}{La pregunta de diseño}
  \begin{alertblock}{Pregunta guía}
    Cuando se comparan tres métodos de empaque, ¿qué convierte una
    diferencia de medias en un efecto de tratamiento y no en una diferencia
    de quién eligió qué?
  \end{alertblock}

  \vspace{0.6em}
  Al final de esta lección, usted puede:
  \begin{itemize}
    \item nombrar el factor, sus niveles, las unidades y la respuesta;
    \item asignar al azar 36 estaciones a tres tratamientos;
    \item leer medias de grupo tras la aleatorización;
    \item explicar por qué los grupos autoseleccionados no son un
      experimento aleatorizado.
  \end{itemize}

  \vfill
  \scriptsize Todas las estaciones, experiencias y tiempos de ciclo usados hoy son sintéticos. No se usan datos reales de empresa.
\end{frame}

\begin{frame}{Un factor de tres niveles y 36 unidades}
  \small
  \begin{center}
    \begin{tabular}{lcc}
      \toprule
      \textbf{Pieza de diseño} & \textbf{En este experimento} & \textbf{Rol} \\
      \midrule
      Factor & Método de empaque & Lo que asignamos \\
      Niveles & Estándar, Guiado, Automatizado & $k=3$ tratamientos \\
      Unidades & 36 estaciones de empaque & A qué asignamos \\
      Respuesta & Tiempo de ciclo (segundos) & Lo que medimos \\
      \bottomrule
    \end{tabular}
  \end{center}

  \vspace{0.5em}
  \begin{exampleblock}{Línea de empaque sintética}
    La asignación aleatoria coloca 12 estaciones en cada método. La
    experiencia del operador es una covariable previa. No se usa para
    elegir el método.
  \end{exampleblock}
\end{frame}

\begin{frame}[fragile]{Paso 1: Asignación aleatoria}
  \begin{columns}[T]
    \begin{column}{0.40\textwidth}
      \small
      \begin{lstlisting}
labels = np.repeat(
  treatments, 12
)
rng.shuffle(labels)
      \end{lstlisting}

      \begin{block}{Salida verificada}
        $n=12$ estaciones por método\\
        Experiencia media $4.618$, $5.943$, $5.666$\\
        Brecha Auto-Estándar $1.047694$ años
      \end{block}
    \end{column}
    \begin{column}{0.60\textwidth}
      \centering
      \includegraphics[width=\textwidth]{../figuras/diseno_experimental_01_factor_aleatorizacion.png}
      \vspace{0.2em}
      \scriptsize Asignación equilibrada y chequeo de covariable del Paso 1
    \end{column}
  \end{columns}
\end{frame}

\begin{frame}{La aleatorización hace de la comparación un contraste de tratamiento}
  \small
  Efectos aditivos verdaderos sobre el tiempo de ciclo:
  \[
    0\ \text{(Estándar)},\quad -2\ \text{(Guiado)},\quad -7\ \text{(Automatizado)}.
  \]
  Tras la asignación aleatoria con semilla 42, las medias observadas son
  $51.871576$, $50.039332$ y $43.852371$ segundos.
  \[
    \bar y_{\text{Automatizado}}-\bar y_{\text{Estándar}}=-8.019206.
  \]
  La brecha está cerca del efecto verdadero Automatizado de $-7$. La
  aleatorización, no un ajuste posterior, es lo que justifica leerla como
  una comparación de método.
\end{frame}

\begin{frame}[fragile]{Paso 2: Tiempos de ciclo aleatorizados}
  \begin{columns}[T]
    \begin{column}{0.40\textwidth}
      \small
      \begin{lstlisting}
y = (
  52 + effect
  - 0.35 * (
    exp - exp.mean()
  ) + noise
)
      \end{lstlisting}

      \begin{block}{Salida verificada}
        Medias $51.872$, $50.039$, $43.852$\\
        Brecha aleatorizada $-8.019206$\\
        Efecto verdadero Automatizado $-7$
      \end{block}
    \end{column}
    \begin{column}{0.60\textwidth}
      \centering
      \includegraphics[width=\textwidth]{../figuras/diseno_experimental_02_resultados_aleatorizados.png}
      \vspace{0.2em}
      \scriptsize Tiempos de ciclo aleatorizados del Paso 2
    \end{column}
  \end{columns}
\end{frame}

\begin{frame}{Paso 3: La autoselección no es un tratamiento}
  \begin{columns}[T]
    \begin{column}{0.42\textwidth}
      \small
      Experiencia media autoseleccionada: $2.678$, $5.655$, $7.895$ años.

      \vspace{0.4em}
      Brecha observacional $Y$ Auto-Estándar: $-14.809630$.

      \vspace{0.4em}
      Brecha de experiencia $5.216727$ frente a aleatorizada $1.047694$.

      \vspace{0.5em}
      \textbf{Límite:} un factor observacional no es un tratamiento
      aleatorizado. La Lección 44 analiza la respuesta aleatorizada de tres
      grupos con ANOVA.
    \end{column}
    \begin{column}{0.58\textwidth}
      \centering
      \includegraphics[width=\textwidth]{../figuras/diseno_experimental_03_limite_autoseleccion.png}
      \vspace{0.2em}
      \scriptsize Medias aleatorizadas frente a autoseleccionadas del Paso 3
    \end{column}
  \end{columns}
\end{frame}

\begin{frame}{Verificación de conceptos}
  \begin{enumerate}
    \item ¿Cuál es el factor experimental y cuáles son sus niveles?
    \item ¿Por qué importa la asignación aleatoria de 36 estaciones para la
      experiencia?
    \item La brecha aleatorizada Auto-Estándar es $-8.019$ segundos y el
      efecto verdadero es $-7$. ¿Por qué pueden diferir?
    \item ¿Por qué la brecha observacional $-14.810$ no es un efecto de
      tratamiento?
  \end{enumerate}

  \vspace{0.6em}
  \begin{block}{Conexión con la Lección 44}
    La sesión puede continuar directamente con ANOVA de un factor, sumas de
    cuadrados y el cociente $F$ para las tres medias aleatorizadas.
  \end{block}

  \vfill
  \scriptsize Fuente: OpenStax, \textit{Introductory Business Statistics 2e},
  Capítulo 12, Sección 12.2. CC BY-NC-SA.
\end{frame}

\appendix



\end{document}
